\begin{figure}[htp!]
  \centering
  \includegraphics[scale=0.22]{Esercizio5.3/Esercizio.jpg} 
\end{figure}
La richiesta di reattore critico si traduce in:
\begin{equation}
    \nu q = 1
\end{equation}
Dove $\nu$ rappresenta il numero di neutroni prodotti per fissione che è circa: $\nu\approx 2.4\,\mathrm{n}/\mathrm{fiss}$. Il $q$ invece è dato da due fattori:
\begin{equation}
    q=(1-p_{\mathrm{catt}})_{\mathrm{mod}}\frac{\Bar{\sigma}_{\mathrm{fiss}}}{\Bar{\sigma}_{\mathrm{fiss}}+\Bar{\sigma}_{\mathrm{n},\gamma}}
\end{equation}
Noi sabbiamo che $P_{\mathrm{catt}}=0.47$, quindi:
\begin{equation*}
    \frac{\Bar{\sigma}_{\mathrm{fiss}}}{\Bar{\sigma}_{\mathrm{fiss}}+\Bar{\sigma}_{\mathrm{n},\gamma}}=0.79
\end{equation*}
Di conseguenza:
\begin{equation}
    0.21\Bar{\sigma}_{\mathrm{fiss}}=0.79\Bar{\sigma}_{\mathrm{n},\gamma}
\end{equation}
Indichiamo con $\alpha$ la frazione di $\ch{^{235}U}$. Quindi:
\begin{equation}
    0.21\alpha\sigma_{\mathrm{fiss},235}=0.79\left(\alpha\sigma_{\mathrm{n},\gamma,235}+(1-\alpha)\sigma_{\mathrm{n},\gamma,238}\right)
\end{equation}
Raccogliendo $\alpha$:
\begin{equation}
    \alpha=4.5\%
\end{equation}
Ora vogliamo sapere a quale tempo $\alpha(t)=0.045$:
\begin{equation}
   \alpha(t)=0.045=\alpha(0)\frac{e^{-\frac{\ln{2}}{\tau_{1/2,235}}t}}{e^{-\frac{\ln{2}}{\tau_{1/2,238}}t}}
\end{equation}
Invertendo la relazione si ottiene:
\begin{equation}
    t=2.25\cdot10^{9}\,\mathrm{yr}
\end{equation}